\subsection{收缩定理的证明}
\label{app:tensor}

在定理~\ref{thm:tensor-contraction} 的证明中，
$\log$ 表示以二为底的对数，$\ln$ 表示自然对数。

\subsubsection{系数环与精确算术}
\label{subsec:tensor-arithmetic}

写 $\lambda=\lceil\log_2 s\rceil$、$e=4\lambda$。
取一个次数为 $e$ 的首一多项式
$f\in\mathbb Z[U]$，其首项系数之下的系数在 $\{0,1\}$
中，且其模二约化不可约。
这可以在 $s$ 的多项式时间内确定性完成：
枚举 $2^e$ 个候选，并用 $\mathbb F_2$ 上每个次数为正、
至多 $e/2$ 的首一多项式检验可除性。
有限域有每个正次数的不可约多项式~\cite{vzGG13}。
置
\[
 \mathcal O=\mathbb Z_{(2)}[U]/(f),\qquad
 \mathcal K=\mathbb Q[U]/(f),\qquad
 \mathcal A=\mathcal O[V]/(V^\ell),\qquad
 \mathcal A_p=\mathcal A/(2^p).
\]
这里 $\mathbb Z_{(2)}$
由分母为奇数的有理数组成。
模二不可约蕴含在 $\mathbb Q$ 上不可约，
故 $\mathcal K$ 是域。
对 $\mathcal K$ 的非零元素，设 $v_2$
为其在 $1,U,\ldots,U^{e-1}$
中坐标的寻常二进赋值的最小值。
$\mathcal O$ 模二的约化是域 $\mathbb F_{2^e}$。
对最小坐标赋值取商因此表明
\[
 v_2(xy)=v_2(x)+v_2(y)\qquad(x,y\ne0).
\]
$\mathcal O$ 中模二非零的每个元素都是单位，
因为其乘法矩阵的行列式在 $\mathbb Z_{(2)}$
中可逆，故逆矩阵的分量在该环中。

$\mathcal A_p$ 中的算术在 $s,p,\ell$
上花费多项式时间与空间。
我们使用基 $1,U,\ldots,U^{e-1}$
与 $1,V,\ldots,V^{\ell-1}$，
并把每个坐标表示为它在零与 $2^p-1$
之间的整数剩余。标量单位的求逆可先模二求逆
再用 Newton 提升完成，或用其乘法矩阵完成，
使用精确的多项式与整数运算~\cite{vzGG13}。
$\mathcal A$ 中元素的可除性与精度总是逐坐标理解。

把 $2^e$ 个剩余域元素的二值坐标提升，
即得两两之差为单位的插值点。
因此，在这些点的任意所需数量上的 Lagrange
插值在 $\mathcal A_p$ 上有效。
我们将使用少于 $s^4$ 个点。
初始分量属于 $\mathcal B_{P_0,\ell}$
的收缩，在 $\mathcal A_{P_0}$ 中计算之后仍属于该子环，
因为它是那些分量的多项式。

精度增大时，我们逐个提升每个张量分量。
对原来按模 $2^p$ 指定的分量，
其\emph{典范提升}是由其全部坐标的整数代表得到的
$\mathcal A$ 中元素。
典范提升不必保持加法或乘法。
每次递归调用按该调用的精度对其指定的张量作收缩。
为定义其变形，它先作出这一逐分量的提升选取。
该调用中后续的全部求值都指称这些相同的提升。
更明确地，每个张量描述都带有其创建精度 $p_c$
（原始张量为 $P_0$）。
在其后的每个精度处，其分量按模 $2^{p_c}$
求值，并使用这些相同的整数代表。
精度沿递归路径严格递增，
故这些代表在其后每个精度处也是典范代表。
特别地，一个包张量在整个递归中保持其创建时的代表。

\subsubsection{参数与碰撞装箱}
\label{subsec:tensor-packing}

下文全部估计对足够大的 $s$ 成立，
阈值只依赖固定常数。
较小规模的固定有限范围用直接求和处理。
参数不等式可以直接检验。置
\[
 H=\left\lceil\frac{s}{\lfloor\sqrt\lambda\rfloor}\right\rceil,
 \qquad r=64,\qquad \varepsilon=\frac1{128\cdot1024},
 \qquad \Delta=\lceil s^{9/10}\rceil.
\]
有 $M\le H$ 条边的调用是叶，直接求值。
在深度 $d$ 的活动调用处（根在深度零），使用
\begin{equation}\label{eq:tensor-parameters}
 u=u_d=\left\lfloor\frac{\varepsilon\lambda}{2^d}\right\rfloor,
 \qquad a=a_d=2\sum_{j=0}^d u_j,\qquad
 a_{-1}=0,\qquad L=2^a,\qquad D_{\rm arr}=r^{2a}.
\end{equation}
我们将安排每个孩子至多 $M/8$ 条边、最大度至多
$\Delta$。这些断言蕴含，在每条活动路径上，
\[
 H<M\le Cs\,8^{-d},\qquad
 d\le\frac{\log_2\lambda}{6}+O_C(1).
\]
因此，对一切足够大的 $s$，
\begin{equation}\label{eq:tensor-parameter-bounds}
 u=\Theta_C(\lambda/2^d),\quad
 \lambda^{4/5}\le u\le M,\quad
 a\le4\varepsilon\lambda,\quad
 L,D_{\rm arr}\le s^{1/1000}.
\end{equation}
特别地，在每次活动调用处
\begin{equation}\label{eq:tensor-packing-conditions}
 M\ge2D_{\rm arr}^{\,2}\Delta,\qquad
 M\ge L D_{\rm arr},\qquad 2L\le\Delta.
\end{equation}
例如，第一个不等式右端是
$O(s^{9/10+2/1000})=o(H)$。
其余不等式同理。

\begin{lemma}[装箱成包]\label{lem:tensor-packets}
在 \eqref{eq:tensor-packing-conditions} 下，
把网络每侧的顶点用 $[r]^a$ 的元素标记。
存在可确定性构造的标记与把边嵌入 $L$
个槽位的包的方案，满足以下性质。
\begin{enumerate}
\item 包至多 $3M/L$ 个。每条边在其行端与列端占据
      同一槽位。未用的槽位没有端口。
\item 在每个包的每个坐标中，行符号与列符号的无序对
      是固定的。当两符号不同时，槽位号的第 $i$
      位指定它们在坐标 $i$ 中的取向
      （$0\le i<a$）。
\item 对每个坐标与每个符号，每侧属于带该符号顶点的
      实端口至多 $3M/r$ 个。
\end{enumerate}
构造在 $s+\mathcal L$ 上使用多项式时间与空间。
\end{lemma}
\begin{proof}
对每对标号，设 $X_\gamma$ 为具有该对的边数。
独立均匀的顶点标号给出
\[
 \mathbb E\sum_\gamma (X_\gamma-M/D_{\rm arr})^2\le2M\Delta.
\]
事实上，把该期望展开为有序边对上的标号一致概率之和
减去 $1/D_{\rm arr}$。端点不交的对贡献零。
共享一个端点的有序对（含相等的边）至多 $2M\Delta$
个，每个贡献至多一。
条件期望给出具有该界的标记。
为实现之，尝试下一个顶点的 $r^a$ 个标号，
并通过枚举边对计算条件期望。
每个概率只涉及至多四个顶点，
故这是精确的多项式时间计算。
这就是通常的条件期望法~\cite{Raghavan88}。

每个负载至多 $M/D_{\rm arr}+\sqrt{2M\Delta}$。
把全部对—标类垫到共同的槽数 $H_{\rm pad}$
（最大负载向上取整到 $L$ 的倍数）。
假设蕴含
\[
 H_{\rm pad}\le M/D_{\rm arr}+\sqrt{2M\Delta}+L
              \le3M/D_{\rm arr}.
\]
固定每个坐标中符号的无序对。
若有 $a'$ 个坐标的符号不等，其取向给出 $2^{a'}$
个对标号。把这些垫后的槽位分装成包：
每包用每个这样的标号取 $2^{a-a'}$ 个槽位。
这是可能的，因为 $H_{\rm pad}$ 被 $L$ 整除。
把一个包的位置从 $0$ 到 $L-1$ 编号。
其第 $i$ 位在第 $i$ 对不等时指定取向，
相等时指定副本。每个包内使用同一取向约定。

包数为 $H_{\rm pad}D_{\rm arr}/L\le3M/L$。
在固定坐标中，给定符号在两侧的全部垫后对标号中恰占
$1/r$。其实端口数因此至多
$H_{\rm pad}D_{\rm arr}/r\le3M/r$。
\end{proof}

\subsubsection{Cauchy 变形与整的局部坐标}
\label{subsec:tensor-deformation}

固定一个活动调用、其包与其逐分量提升的张量。
使用
\[
 S=\{-L,-L+1,\ldots,L-1\},\qquad
 \mathcal C(T)=\prod_{t\in S}(T-t),\qquad \alpha=2^{a+1}.
\]
在一个包内，对其行槽位 $x$ 与列槽位 $y$，置
\begin{equation}\label{eq:tensor-cauchy}
 W_{xy}(T)=\frac{\mathcal C(T)}
                  {\mathcal C'(0)(T-(y-x))}.
\end{equation}
分母中的因子消去分子中的一个因子，因为 $y-x\in S$。
故这是 $\mathbb Q$ 上次数 $2L-1$ 的多项式。
把不同包之间的分量置零，并限制到实端口。
在原来的边等同关系下，$W(0)=I$：
对角分量为二，其他每个分量都有一个未消去的零因子。

设 $A(Y)$ 与 $B(X)$ 分别是提升后的行张量与列张量
在独立端口赋值 $Y,X\in\{0,1\}^M$ 上的乘积。
写 $l_0(z)=1-z$、$l_1(z)=z$。定义
\begin{equation}\label{eq:tensor-deformed-contraction}
 F(T)=\sum_{X,Y\in\{0,1\}^M} A(Y)B(X)
                  \prod_{\text{行端口 }x}l_{Y_x}((W(T)X)_x).
\end{equation}
其系数落在 $\mathcal K[V]/(V^\ell)$ 中，
次数至多 $M(2L-1)$，而 $F(0)$
是提升后张量的收缩。特别地，
把 $F(0)$ 模该调用的精度化简即得所需答案。

\begin{lemma}\label{lem:tensor-integrality}
对每个 $t\in S$，$W(t+\alpha z)$ 的分量属于
$\mathcal O[z]$，且 $F(t+\alpha z)$ 属于
$\mathcal A[z]$。
给定 $\zeta$ 模 $2^p$ 的 $e$ 个坐标，
每个分量 $W_{xy}(t+\alpha\zeta)\bmod 2^p$
可以用 $s,p$ 上多项式次位运算计算。
\end{lemma}
\begin{proof}
网格 $S$ 是模 $2^{a+1}$ 的完全剩余系。
对 $S$ 中 $t\ne q$，$v_2(t-q)\le a$，
且每个赋值处非零差的个数与 $t$ 无关。
故 $v_2(\mathcal C'(t))=v_2(\mathcal C'(0))$。
在
\[
 \frac{\mathcal C(t+\alpha z)}{\mathcal C'(0)}
  =\frac{\alpha z\prod_{q\in S\setminus\{t\}}
                         (t-q+\alpha z)}{\mathcal C'(0)}
\]
中，每个以 $q\ne t$ 为下标的因子逐系数被
$2^{v_2(t-q)}$ 整除。对这些幂与分母相消后，
所示多项式被 $2^{a+1}$ 整除。
为得 \eqref{eq:tensor-cauchy}，只省略以 $y-x$
为下标的因子。若 $y-x=t$，这去掉 $\alpha z$
并留下整系数多项式。否则 $\alpha z$
保留，而被省略因子有二进容量
$2^{v_2(t-y+x)}$（$v_2(t-y+x)\le a$），
故结果仍被二整除。代入
\eqref{eq:tensor-deformed-contraction}
即证明关于 $F$ 的断言。

对求值，先把每个以 $q\ne t$ 为下标的因子
$t-q+\alpha z$ 除以已知整数 $2^{v_2(t-q)}$。
所得一次多项式的两个系数都是整数。
如上对分母消去其余二的显式幂，
并把其奇部模 $2^p$ 求逆。
现在可以直接在所给 $\zeta$
的剩余处求值该乘积（含 $\zeta=0$）。
输入有 $ep$ 位，乘积有 $O(L)$ 个因子。
全部坐标运算在 $s,p$ 上花费多项式时间。
\end{proof}

\subsubsection{锥形 Hermite 恢复}
\label{subsec:tensor-hermite}

下一条引理单列出精度估计。
其陈述中全部幂级数截断都是 $z$
中的通常截断。

\begin{lemma}[锥形 Hermite 恢复]\label{lem:tapered-hermite}
设整数 $u,M$ 满足 $1\le u\le a$、$M\ge u$，
并置 $h=\lceil3M/u\rceil$。
对 $t\in S\setminus\{0\}$ 置
\[
 i=v_2(t),\qquad q_t=q(i)=h\max(0,u-i),\qquad D_t=L/2^i.
\]
设 $F\in(\mathcal K[V]/(V^\ell))[T]$
次数至多 $M(2L-1)$，$F(0)\in\mathcal A$，
且对每个 $t\in S$ 有
$F(t+\alpha z)\in\mathcal A[z]$。
对 $p\ge1$，其值 $F(0)$ 模 $2^p$
可以从
\[
 F(t+\alpha z)\pmod{2^{p_t}},
 \qquad
 p_t=p+8hD_t+8M(a+1),\qquad q_t>0
\]
恢复。只需每个多项式的前 $q_t$ 个系数。
恢复在 $M,L,\ell,e$ 与各精度上花费多项式时间与空间。
\end{lemma}
\begin{proof}
$S$ 中有 $L$ 个奇点，故
$\sum_{t\ne0}q_t\ge Lhu\ge3ML>M(2L-1)$。
且 $h\le4M/u$，故 $q_t\le4M$。
对每个正重数的 $t$ 置
\[
 E_t(T)=\prod_{q\in S\setminus\{0,t\}}
                    \left(\frac{T-q}{t-q}\right)^{q_q}.
\]
Hermite 插值给出
\begin{equation}\label{eq:tensor-hermite-origin}
 F(0)=\sum_{t:q_t>0}E_t(0)
       \left(\left[\frac{F(t+\alpha z)}
                            {E_t(t+\alpha z)}\right]_{<q_t}\right)
                            (-t/\alpha).
\end{equation}
为完整起见，把右端的零换成未定元 $T$，
把 $-t/\alpha$ 换成 $(T-t)/\alpha$。
每个被加项次数小于 $\sum q_t$；
在其自身节点处有指定的 jet，
在其他每个节点处按所需阶为零。
它与 $F$ 之差被
$\prod_{t\ne0}(T-t)^{q_t}$ 整除且次数更小，
故为零。这是 Hermite 插值的通常中国剩余构造
~\cite[Section~5.6]{vzGG13}。

$E_t(t+\alpha z)$ 的每个因子是
$1+\alpha z/(t-q)$ 的幂。它是整的且常项为一，
其截断到 $q_t$ 的倒数也是整的，
故作成 \eqref{eq:tensor-hermite-origin}
中的商不损失精度。

固定 $i=v_2(t)$ 与 $D=D_t$。直接从乘积，
\begin{align}
 v_2(E_t(0))
 &=\sum_{q\ne0,t}q_q\bigl(v_2(q)-v_2(t-q)\bigr)\notag\\
 &=D\sum_{j=1}^{a-i}\frac{j}{2^j}
                           \bigl(q(i+j)-q(i)\bigr)
 \ \ge -6Dh. \label{eq:tensor-hermite-loss}
\end{align}
第二等式的理由：赋值小于 $i$ 的点贡献零。
对 $1\le j\le a-i$，有 $D/2^j$
个 $v_2(q)=i+j$ 的点，同样个数的
$v_2(t-q)=i+j$ 的点。后者有 $v_2(q)=i$，
而这些层穷尽了该赋值处除 $t$
外的全部点。这些计数来自完全剩余网格：
其非零整数差的赋值至多 $a$。
最后，$q(i+j)-q(i)\ge-hj$ 且
$\sum_{j\ge1}j^2/2^j=6$。

次数低于 $q_t$ 的多项式在 $-t/\alpha$
处求值至多再损失 $4M(a+1)$ 个二的幂。
于是把每个商系数换成其模 $2^{p_t}$
剩余的任意提升，只会把其在
\eqref{eq:tensor-hermite-origin}
中的贡献改变一个 $2^p\mathcal A$
中的元素。对这些误差求和不损失更多精度。
重构的和是整的，因为它与整值 $F(0)$
相差一个 $2^p\mathcal A$ 中的元素。

最后这一组合可以用所指定基下的有理算术执行。
乘子 $E_t(0)(-t/\alpha)^j$（$j<q_t$）
是计重数的 $O(a)$ 位有理因子的多项式个乘积。
其分子与分母长度因此在 $M,L,a$
上是多项式的。这样的系数只有多项式个，
故即使其分母的乘积也是足够的多项式位公分母。
相加之后，消去其二进幂，
并把所得整坐标模 $2^p$ 化简。
这证明了算法断言。
\end{proof}

在输入精度为 $p$ 的调用处，
以 \eqref{eq:tensor-parameters} 的参数应用该引理。
对每个活动的 $t$，
通过在剩余域元素的 $M(2L-1)+1$
个互异的二值坐标提升处求值再做 Lagrange 插值，
得到整个多项式 $F(t+\alpha z)$ 模 $2^{p_t}$。
引理~\ref{lem:tensor-integrality}
容许该插值在 $\mathcal A_{p_t}$ 上进行。
由 \eqref{eq:tensor-parameter-bounds}，
点数小于 $s^4$。
下文的求值算法把精度 $p_t$ 传给其孩子。

输入精度与孩子精度满足
\begin{equation}\label{eq:tensor-precision-bound}
 p\le s\lambda^{O(1)}2^{a_{d-1}},
 \qquad
 p_t=O_C(MD_t/u).
\end{equation}
第一个断言在根处成立，因为 $P_0\le Cs$。
沿一条路径，精度增量是 $8hD_t+8M(a+1)$，
其中 $h\le4M/u$、$M\le Cs$、$D_t\le2^{a_d}$。
在 $O(\log\lambda)$ 层深度上求和，
即对每个调用（含叶）证明第一个断言。
对第二个：$q_t>0$ 蕴含 $i<u$，故
\[
 D_t\ge2^{a_d-u}=2^{a_{d-1}+u}.
\]
由 $M>H$ 与 $u\ge\lambda^{4/5}$，
这一额外的因子 $2^u$
支配 $\lambda$ 的每个固定幂。
于是 $p$ 的第一个界以及 $8M(a+1)$
都是 $O_C(MD_t/u)$。
余下的增量 $8hD_t$ 直接具有该界。
由 $a_d\le4\varepsilon\lambda$，
全部精度都保持在 $s$ 的多项式内。

锥形化解释了该估计中的因子 $1/u$。
作为对比，给 $u$ 以下的每一层指派同一重数
$q_*=\Theta(M)$，其余层指派零。
在最后保留的层 $i=u-1$ 处，
恒等式 \eqref{eq:tensor-hermite-loss}
将给出损失
\[
 -q_*D_t\sum_{j=1}^{a-i}\frac{j}{2^j}=-\Theta(MD_t).
\]
而我们的重数每层只变化 $h=\Theta(M/u)$，
给出损失 $O(MD_t/u)$。
$a_d$ 的累积选取提供了吸收继承精度与求值损失所需的
额外因子 $2^u$。
因此在每个保留层、每个递归深度处都成立较小的界。

\subsubsection{组间相互作用的模}
\label{subsec:tensor-modes}

固定活动 $t$ 与 $z$ 的插值取值
$\zeta\in\mathcal O$。
本小节中 $T=t+\alpha\zeta$，
全部计算按模 $2^{p_t}$ 进行。
置 $i=v_2(t)<u$、$D=L/2^i$。
两侧顶点按其在坐标 $i$ 中的符号分组。
$W(T)$ 中两端口同组的分量称为局部的，
保持不变。

这一坐标选取使 $W(t)$ 是局部的。
事实上，除非 $y-x=t$，否则 $W_{xy}(t)=0$。
加一个赋值为 $i$ 的整数会翻转第 $i$ 位，
故存留的对有不同的第 $i$ 位。
当包在该坐标中的两个符号不同时，
这意味着行符号与列符号一致。
符号相等的对本来就同组。
从 $t$ 走到 $t+\alpha\zeta$
引入我们现在展开的交叉项。

在其坐标 $i$ 中两符号不同的包里，
行位置 $x$ 与列位置 $y$ 不同组
当且仅当它们的第 $i$ 位相等。
把列分划成瓦片（tile）
\[
 y=y_0+2^{i+1}j,\qquad 0\le y_0<2^{i+1}.
\]
列组在瓦片上为常数。对跨越到该瓦片的行，
\[
 v_2(t-y_0+x)\le i.
\]
否则 $y_0$ 将与 $x+2^i$ 模 $2^{i+1}$
一致，与其第 $i$ 位相等矛盾。
加 $\alpha\zeta$ 不改变此赋值。
几何展开因此给出
\begin{equation}\label{eq:tensor-mode-expansion}
 W_{xy}(T)=
 \sum_{\substack{b\ge0\\b+1<p_t}}
 \frac{\mathcal C(T)}
      {\mathcal C'(0)(T-y_0+x)}
 \left(\frac{2^{i+1}}{T-y_0+x}\right)^b j^b
 \pmod {2^{p_t}}.
\end{equation}
事实上，第一个分式的赋值至少 $a+1-i\ge1$，
而该比值取 $b$ 次幂前赋值至少为一。
几何余项被 $2^{p_t}$ 整除。
这些表达式可以经分解已知的二的幂、
只用单位逆而计算。

对每个这样的包、瓦片与 $b$，
引入权重 $w_Z=b+1$ 的模变量 $Z$，
其将替换为
\[
 L_Z(X)=\sum_{\text{瓦片中的实列 }y}j^bX_y.
\]
照常 $j^0=1$（含 $j=0$）。
其源是该瓦片的列组，其目标是该包坐标对中的另一组。
公式 \eqref{eq:tensor-mode-expansion}
把每个行型表示为其局部部分加上其目标模的线性组合。
$Z$ 的系数被 $2^{w_Z}$ 整除。
任一权重的变量至多
\begin{equation}\label{eq:tensor-mode-multiplicity}
 (3M/L)2^{i+1}=6M/D
\end{equation}
个。

在仿射因子 $l_{Y_x}((W(T)X)_x)$
的乘积中，模单项 $Z^\beta$
的每个系数都被
$2^{\sum_Zw_Z\beta_Z}$ 整除。
因此只需考虑满足
\begin{equation}\label{eq:tensor-mode-budget}
 \sum_Zw_Z\beta_Z<p_t
\end{equation}
的指数元组。这些元组的个数是
\begin{equation}\label{eq:tensor-partition-count}
 \exp\bigl(O(\sqrt{Mp_t/D})\bigr)
       =\exp\bigl(O_C(M/\sqrt u)\bigr).
\end{equation}
以下是完整的计数估计。对任何 $\theta>0$，
用生成乘积与 \eqref{eq:tensor-mode-multiplicity}
把个数界定为
\[
 \exp\left(\theta p_t+\frac{6M}{D}
                   \sum_{j\ge1}-\ln(1-e^{-\theta j})\right).
\]
展开对数并对几何级数求和，
\[
 \sum_{j\ge1}-\ln(1-e^{-\theta j})
  =\sum_{h\ge1}\frac1{h(e^{\theta h}-1)}
  \le\frac1\theta\sum_{h\ge1}\frac1{h^2}.
\]
取 $\theta=\sqrt{M/(Dp_t)}$ 并用
\eqref{eq:tensor-precision-bound}。
全部模变量及其指数都有多项式规模描述。
满足 \eqref{eq:tensor-mode-budget}
的元组可以按多项式延迟与空间枚举：
对变量排序，逐个在剩余预算内选取每个指数。
每个可行的部分选取都有全零的延续，
故被访问的部分选取数至多为完整元组数乘以变量数
（差一个可加常数）。

\subsubsection{孩子网络及其描述}
\label{subsec:tensor-children}

对 \eqref{eq:tensor-mode-budget} 中的一个元组
$\beta$，为每个组 $g\in[r]$
构造一个孩子网络。把当前网络两侧属于组 $g$
的全部当前顶点放到孩子的同一侧。
保留其端口与所选提升分量（现按模 $2^{p_t}$
化简）。另一侧为每个包放一个顶点，
连边到该包的每个保留端口。
其下标恰为该包中组 $g$ 的行位 $Y$
与列位 $X$。

为定义这个新的包张量，把包中组 $g$
的行因子 $l_{Y_x}$ 相乘，
其自变量用局部列位与形式模表示。
提取以该包中所有目标为 $g$ 的模的指数 $\beta_Z$
的系数。把它乘以
\[
 \prod_{\substack{Z\text{ 在此包中}\\
                  \text{source}(Z)=g}}L_Z(X)^{\beta_Z}.
\]
空乘积为一。若无目标行，则只要某目标指数为正，
提取的系数为零。若 $Z$
的源瓦片没有实列，则 $L_Z(X)=0$，
其幂在 $\beta_Z=0$ 时贡献一，否则贡献零。

每个模恰有一个目标与一个源。
因此系数提取在「包—组」行乘积上可分解，
且每个模的代入只依赖其源组中的位。
在 \eqref{eq:tensor-deformed-contraction}
中展开乘积，随后在组内对独立位求和，
证明精确恒等式
\begin{equation}\label{eq:tensor-recursion-identity}
 F(t+\alpha\zeta)
   =\sum_{\beta:\,\sum w_Z\beta_Z<p_t}
                       \prod_{g=1}^r Z_{g,\beta}
                  \pmod {2^{p_t}},
\end{equation}
其中 $Z_{g,\beta}$ 是刚才定义的孩子的收缩。
被省略的单项由其系数可除性而为零。

由引理~\ref{lem:tensor-packets}，
孩子至多 $6M/r=3M/32<M/8$ 条边。
旧顶点的度不变，新的包顶点至多 $2L\le\Delta$
个端口。于是 \eqref{eq:tensor-parameter-bounds}
所用的不变量归纳成立。

全部新张量都是隐式给出的。
为求值一个包张量分量，展开其至多 $L$
个仿射行因子的乘积。每个因子在 $s$
上有多项式个项（含其局部项）。
枚举这些选取、检验所需指数并对所得项求和花费
\[
 \exp(O(L\log s))
\]
时间与多项式空间。因子 $L_Z(X)^{\beta_Z}$
用 $\mathcal A_{p_t}$ 中通常的反复平方计算。
由 $L\le s^{1/1000}$，该分量成本为
$\exp(O(s^{9/10}))$。

新张量的描述由其包数据、元组 $\beta$、
插值取值与创建精度组成。
这些数据与局部核足以直接计算其分量。
保留的张量从其原描述求值，
或从其自己的创建时描述求值（若它在递归中较早创建）。
其在那个精度处的剩余随后被典范提升到当前精度。
旧描述不可变并按引用共享。
因此分量例程的复合不随深度增长，
全部分量例程保持所需的时间与空间界。
沿一条路径只添加多项式个顶点与多项式多的新描述数据。

\subsubsection{时间与空间界}
\label{subsec:tensor-cost}

\begin{proof}[定理~\ref{thm:tensor-contraction} 的证明]
把原始分量嵌入 $\mathcal A_{P_0}$，
执行上述递归算法。在活动调用处，
对每个正重数的节点与每个所需的、剩余互异的插值取值，
用其孩子收缩计算
\eqref{eq:tensor-recursion-identity}。
插值 $F(t+\alpha z)$ 并应用
引理~\ref{lem:tapered-hermite}。
在叶处，直接对全部布尔边赋值求和。
引理~\ref{lem:tensor-integrality}、
递归恒等式与恢复引理由叶向上归纳证明正确性。

一次调用处的插值节点数、取值数与组数是 $s$
的多项式，与 $\mathcal L$ 无关。
由 \eqref{eq:tensor-partition-count}，
深度 $d$ 的活动调用的孩子数至多
\[
 s^{O(1)}\exp\left(
       O_C\left(\frac{s}{\sqrt\lambda}
                    \left(\frac{\sqrt2}{8}\right)^d\right)\right).
\]
事实上 $M\le Cs\,8^{-d}$
且 $u_d=\Theta_C(\lambda/2^d)$。
把这些界在 $O_C(\log\lambda)$ 层深度上相乘，
或以此方式界定每层节点数，得总调用数至多
\[
 \exp\left(O_C(s/\sqrt\lambda)
                   +O_C(\lambda\log\lambda)\right)
       =\exp(O_C(s/\sqrt\lambda)).
\]
在叶处，直接求和至多用 $2^H$
个赋值，每赋值多项式次分量求值。
其成本被 $s+\mathcal L$
的一个多项式乘以
\[
 \exp(O_C(H)+O_{C,A}(s^{9/10}))
   =\exp(O_{C,A}(s/\sqrt\lambda))
\]
界定。局部算术、装箱、描述与恢复对每个生成的项
花费多项式时间（已算界的分量求值除外）。
依赖描述的因子界定每次调用的工作。
由于分支界与 $\mathcal L$ 无关，
总工作保持 $s+\mathcal L$
的固定多项式因子。这证明时间界。

对空间，深度优先遍历递归，
一次累加一个关于 $\beta$ 的和与一个关于组的积。
存储一次调用的恢复所需的多项式个 Taylor 数据、
其当前插值与模下标、其图与新张量描述。
其规模在 $s+\mathcal L$ 上是多项式的
（精度由 \eqref{eq:tensor-precision-bound}
计入）。栈层只有 $O_C(\log\lambda)$ 层。
引理~\ref{lem:tapered-hermite}
中的公分母计算也只用多项式位整数。
每次调用按其自身的固定精度累加其收缩项。
原始分量例程按假设使用多项式空间，
新例程亦然。这证明空间界。
最后提取 $\mathcal A_{P_0}$
中的基坐标，它们给出 $\mathcal B_{P_0,\ell}$
中的答案。
\end{proof}
